\chapter{核心概念与理论基础}

\section{核心概念界定}

\subsection{泉州世界遗产}

世界遗产是指经联合国教科文组织世界遗产委员会确认，具有突出普遍价值的文物、建筑群或遗址，以及具有突出普遍价值的自然区域。泉州于2021年列入《世界遗产名录》的项目全称为“泉州：宋元中国的世界海洋商贸中心”，由二十二处遗产点构成，包括九日山祈风石刻、市舶司遗址、德济门遗址、天后宫、真武庙、南外宗正司遗址、泉州府文庙、开元寺、老君岩造像、清净寺、伊斯兰教圣墓、草庵摩尼光佛造像、磁灶窑址、德化窑址、安溪青阳下草埔冶铁遗址、洛阳桥、安平桥、顺济桥遗址、江口码头、石湖码头、六胜塔、万寿塔等。这些遗产点从制度、交通、生产、信仰与社群等不同侧面，呈现了宋元泉州作为世界海洋商贸中心的整体格局\cite{quanzhou2021}。

从语文课程的角度看，泉州世界遗产至少具有三重教育价值。第一是历史叙事价值，遗产点背后关联着海洋贸易、移民迁徙、宗教共存与城市发展等宏大叙事，为学生的阅读与探究提供了丰富材料。第二是语言文化价值，楹联、碑刻、地名、传说、闽南方言与民间故事共同构成了多层次的文本世界，学生可以在真实语料中感受语言的多样性与表现力。第三是审美体验价值，石塔、古桥、红砖厝、蚵壳厝等物质遗存与南音、梨园戏、簪花围等非物质文化遗产相互映衬，为学生提供了可看、可听、可写的审美对象。

\subsection{语文综合性学习}

语文综合性学习是语文课程中以任务为载体，整合听说读写、沟通课堂内外、联系书本学习与生活实践的学习领域。它强调在真实或接近真实的情境中，围绕具有一定挑战性的问题或主题，引导学生综合运用语文知识与能力，开展自主、合作、探究的学习活动\cite{wangning2016}。与单篇阅读教学相比，综合性学习的目标更综合、内容更开放、过程更强调实践，评价也更关注学生在完成任务过程中的表现。

《义务教育语文课程标准（2022年版）》将综合性学习纳入“实用性阅读与交流”“文学阅读与创意表达”“思辨性阅读与表达”等学习任务群的整体框架，强调学习内容的整合性与学习方式的实践性\cite{jiaoyubu2022}。这意味着综合性学习不是语文课程的“附加活动”，而是学生形成语文核心素养的重要途径。有效的综合性学习应当具有清晰的语文目标、真实的学习情境、结构化的任务链条和可见的学习证据。

\subsection{项目化学习}

项目化学习是以真实问题或产品为驱动，学生在持续探究、合作交流与反思修正中建构知识、发展能力的学习方式。其核心要素通常包括：指向学科核心概念的驱动性问题、需要持续探究的真实任务、学生的自主选择与协作、可公开呈现的成果以及贯穿全过程的表现性评价\cite{xiaxuemei2018}。项目化学习与一般主题活动的重要区别在于，它不是先学知识再“应用”于活动，而是让知识在解决问题的过程中被需要、被使用、被反思。

在语文学科中，项目化学习需要特别处理“项目逻辑”与“语文逻辑”的关系。项目逻辑关注任务是否真实、成果是否有受众、过程是否可探究；语文逻辑关注学生是否在阅读、表达、交流与反思中提升了语言文字运用能力。二者并非相互排斥，而是需要在设计中有意识地建立联结。例如，学生为校园展览撰写解说词，既是在完成一个真实的传播任务，也是在练习说明性语言的准确性、对象意识与结构安排。

\section{理论基础}

\subsection{文化回应性教学}

文化回应性教学强调教师应关注学生的文化背景、生活经验与社区资源，并将这些内容纳入课程与教学，使学习内容与学生产生意义联结\cite{gay2010}。该理论认为，文化不仅是教学内容的一部分，也是学生理解世界、表达自我的重要资源。对于地方世界遗产教育而言，文化回应性教学提示我们：不能把遗产当作远离学生生活的“客观知识”，而应引导学生从自己的生活经验出发，与遗产建立情感与认知的联系。

在本研究中，文化回应性教学主要体现为三个设计取向。其一，从学生熟悉的地方经验入手，例如从“你走过洛阳桥吗”“家里是否讲过浔埔女的故事”等问题出发，激活学生的先前经验。其二，尊重学生对遗产的多元理解，允许学生提出疑问、表达不同看法，而不是要求他们接受单一的标准答案。其三，将社区资源引入课堂，邀请家长、地方文化研究者或非遗传承人参与交流，使学习发生在学校与社区之间。

\subsection{建构主义学习理论}

建构主义认为，学习不是知识的被动接收，而是学习者在已有经验基础上主动建构意义的过程。维果茨基特别强调社会互动与“最近发展区”的作用，认为学习者在成人或能力更强的同伴支持下，能够完成独立状态下难以完成的任务\cite{vygotsky1978}。这一观点为项目化学习中的支架设计提供了理论依据：教师不应直接给出答案，而应通过问题引导、示范、工具支持和同伴协作，帮助学生逐步承担更复杂的学习任务。

在本研究的项目设计中，建构主义主要影响三个方面。第一，任务设计遵循由易到难、由局部到整体的顺序，让学生在完成子任务的过程中逐步形成对项目的整体理解。第二，提供阅读材料、访谈提纲、观察记录表、写作支架等工具，帮助学生在“最近发展区”内完成挑战。第三，强调同伴互评与小组协商，使学生在交流中修正自己的理解，形成更完整的表达。

\subsection{逆向教学设计}

逆向教学设计主张从预期的学习结果出发，先确定学生应当知道什么、能够做什么，再设计证明学习发生的评价证据，最后规划学习活动\cite{wiggins2005}。这一设计思路与项目化学习具有天然的契合性：项目成果不是教学的终点，而是学生展示理解与能力的证据；教学活动则应围绕这些证据展开，而不是先安排活动再寻找目标。

本研究借鉴逆向教学设计，将项目设计分为三个阶段。第一阶段明确预期结果，即学生应理解泉州世界遗产的文化价值，能够运用说明、描写、议论等表达方式完成讲解与推介，并在合作探究中形成证据意识。第二阶段确定评价证据，包括资料卡片、访谈记录、解说词、导览手册、展示汇报与反思日志。第三阶段规划学习活动，将阅读、走访、访谈、写作、修改与展示嵌入项目流程，使活动始终服务于学习目标。

\section{项目化学习与地方遗产教育的适切性}

从课程目标看，地方世界遗产能够为语文综合性学习提供真实、丰富且具有文化深度的内容。学生围绕遗产点展开阅读与探究，需要在不同文本之间建立联系，需要在事实与观点之间作出区分，也需要用准确、得体的语言向他人介绍自己的发现。这些活动天然包含信息提取、整合解释、表达交流与反思评价等语文实践。

从学习方式看，项目化学习能够把分散的活动组织成具有内在逻辑的任务链条。传统遗产教育常常以讲座、参观或征文为主，活动之间缺乏连续性与递进性；项目化学习则通过驱动性问题、子任务与成果展示，将学习过程串联起来，使学生始终知道“为什么学”“学到什么程度”“如何证明自己学会了”。这种任务链条有助于学生形成整体理解，也有助于教师观察学生的学习过程。

从资源条件看，泉州世界遗产分布广泛，学校周边往往就有可用的遗产点或文化场馆，开展实地走访与访谈具有一定的可行性。与此同时，项目化学习对课时、交通、安全与家校协同提出了更高要求，需要学校在课程安排与资源协调上提供支持。因此，本研究在设计中有意控制项目规模，将实地探究与课堂学习结合，避免因活动过多而挤压语文学习时间。
